\section{Introduction}

Thermoelastic wave propagation in geological media is relevant to geophysics, engineering geology, geothermal studies, and subsurface monitoring. Rocks are not ideal homogeneous solids: their response depends on mineral composition, density, porosity, fractures, bedding, thermal properties, elastic stiffness, and in situ stress-temperature conditions. When a temperature disturbance appears in such a medium, it may produce thermal expansion, thermal strain, stress redistribution, and mechanical motion. At the same time, elastic waves are controlled by density and elastic moduli, while heat transfer is controlled by thermal conductivity and heat capacity. Therefore, the prediction of thermoelastic response requires a formulation that connects geological material properties with thermal and mechanical fields \cite{mavko2009,jaeger2007,aki2002,nowacki1986}.

The relevance of this work is strengthened by the directional nature of geological media. Bedding, foliation, oriented pores, cracks, joints, and lithological contrasts can make wave propagation and heat transfer depend on direction. Even when the mathematical model uses an effective locally homogeneous and isotropic approximation, the source--probe direction remains important for interpreting the response and comparing different prediction routes. Directional prediction is therefore useful as a compact way to study how a thermoelastic disturbance may be observed between selected points in a rock domain.

The research problem of this thesis is the development of an AI-assisted framework for directional prediction of thermoelastic wave propagation in geological media. The work does not aim to replace classical numerical solvers or to produce a field-validated industrial simulator. Instead, it uses the physical theory of thermoelasticity and effective geological material parameters to define a consistent input-output prediction task. Within this task, several neural surrogate approaches can be compared under the same physical scenario, material catalog, source--probe geometry, and output interpretation.

The object of the study is thermoelastic wave propagation in geological media represented by effective material parameters. The subject of the study is the AI-based directional prediction of temperature-related and displacement-related response quantities for selected geological media and source--probe configurations. The thesis focuses on representative rock types such as sandstone, basalt, granite, and limestone, while recognizing that real geological formations are heterogeneous, porous, fractured, layered, and often anisotropic.

The aim of the thesis is to design and analyze a prototype framework for AI directional prediction of thermoelastic wave propagation in geological media. To achieve this aim, the following tasks are addressed: to review the physical and geological factors controlling thermoelastic wave propagation in rocks; to formulate the mathematical variables, assumptions, material parameters, governing relations, and directional quantities of the problem; to define a common prediction methodology and input-output structure; to generate and describe a controlled COMSOL-based simulation dataset; to implement a service-oriented prototype with a unified prediction interface; to compare PINN, FNO, MeshGraphNet, and Transformer-style predictive components; and to evaluate their behavior using travel time, displacement, temperature response, directional prediction, latency, and limitations.

The scientific significance of the work lies in combining geological interpretation, thermoelastic theory, and machine-learning prediction into one consistent framework. The practical significance lies in the implementation of a prototype system that allows several predictive components to be evaluated through one interface. In this comparison, the PINN is treated as the explicitly physics-informed baseline because it includes thermoelastic residuals derived from the governing equations, while the FNO, MeshGraphNet, and Transformer-style components are treated as data-driven or architecture-based baselines that use physically meaningful inputs and outputs without enforcing the full thermoelastic PDE system in the same way.

The structure of the thesis follows this logic. Chapter 2 presents the physical and geological foundations of thermoelastic wave propagation. Chapter 3 introduces the mathematical formulation of the problem. Chapter 4 defines the prediction methodology and model comparison framework. Chapter 5 describes the COMSOL simulation setup and dataset processing. Chapter 6 presents the implementation architecture and prediction interface. Chapters 7--10 describe the PINN, FNO, MeshGraphNet, and Transformer-style components. Chapter 11 presents the web implementation, and Chapter 12 discusses the experimental results. The conclusion summarizes the findings, limitations, and directions for future work.
